\documentclass[11pt]{article}

% ------------------------------------------------------------
% Packages
% ------------------------------------------------------------
\usepackage[margin=1in]{geometry}
\usepackage{amsmath,amssymb}
\usepackage{array}
\usepackage{booktabs}
\usepackage{enumitem}
\usepackage[T1]{fontenc}
\usepackage{lmodern}
\usepackage{microtype}

% ------------------------------------------------------------
% General Formatting
% ------------------------------------------------------------
\setlength{\parindent}{0pt}
\setlength{\parskip}{0.75em}
\setlist[itemize]{leftmargin=1.5em}
\setlist[enumerate]{leftmargin=1.75em}

% ------------------------------------------------------------
% Whiteboard Environment
% Everything placed inside this environment is intended
% to be physically written on the whiteboard during lecture.
% ------------------------------------------------------------
\newenvironment{whiteboard}
{
    \par\medskip
    \begin{center}
    \begin{minipage}{0.94\textwidth}
    \hrule
    \vspace{0.4em}
    \textbf{\large WHITEBOARD}
    \vspace{0.3em}
    \hrule
    \vspace{0.7em}
}
{
    \vspace{0.5em}
    \hrule
    \end{minipage}
    \end{center}
    \medskip
}

% ------------------------------------------------------------
% Document
% ------------------------------------------------------------
\begin{document}

\begin{center}
    {\LARGE \textbf{Computer System Architecture}}\\[0.5em]
    {\LARGE \textbf{Lecture 1}}\\[0.5em]
    {\Large \textbf{Chapter 3: Data Representation}}\\[0.5em]
    \textit{Based on M. Morris Mano, Computer System Architecture, 3rd ed.}
\end{center}

\vspace{1em}

% ============================================================
\section{The Central Problem: What Do Bits Mean?}
% ============================================================

A digital computer ultimately stores information in registers and memory as
binary values. A register consists of binary storage elements, so physically it
contains only patterns of \(0\)'s and \(1\)'s.

Those bit patterns may represent very different kinds of information.

They may represent:

\begin{itemize}
    \item numbers used in arithmetic,
    \item letters and other characters,
    \item discrete symbols,
    \item or control information used to direct computer operations.
\end{itemize}

The bits alone therefore do not tell us what the stored information means.
The meaning depends upon the \emph{representation} being used.

\begin{whiteboard}
\[
\boxed{\text{DATA REPRESENTATION}}
\]

\[
\text{Registers store bits: }0,1
\]

\[
\text{Bits may represent numbers, characters, or other symbols.}
\]

\[
\boxed{\text{bit pattern}+\text{interpretation}=\text{meaning}}
\]
\end{whiteboard}

This idea will run throughout the chapter. The same physical bit pattern can
be interpreted differently depending upon what type of data the computer
expects it to represent.

% ============================================================
\section{Number Systems and Radix}
% ============================================================

A positional number system has a \emph{base}, also called a \emph{radix},
denoted by \(r\).

A number written in radix \(r\) consists of digits whose positions are weighted
by integer powers of \(r\). In general,

\[
N=\sum_i d_i r^i,
\]

where \(d_i\) is the digit occupying position \(i\).

The four number systems of greatest interest here are:

\[
\begin{array}{c|c}
\text{Number System} & \text{Radix} \\ \hline
\text{Binary} & 2\\
\text{Octal} & 8\\
\text{Decimal} & 10\\
\text{Hexadecimal} & 16
\end{array}
\]

In hexadecimal, the symbols

\[
A,B,C,D,E,F
\]

represent decimal values

\[
10,11,12,13,14,15.
\]

\begin{whiteboard}
\[
\boxed{\text{Radix }r}
\]

\[
N=\sum_i d_i r^i
\]

\[
\begin{array}{c|c}
\text{System} & r \\ \hline
\text{binary} & 2\\
\text{octal} & 8\\
\text{decimal} & 10\\
\text{hexadecimal} & 16
\end{array}
\]

\[
A=10,\quad B=11,\quad \ldots,\quad F=15
\]
\end{whiteboard}

For example,

\[
(101101)_2
\]

has the decimal value

\[
1(2^5)+0(2^4)+1(2^3)+1(2^2)+0(2^1)+1(2^0).
\]

Therefore,

\[
(101101)_2=(45)_{10}.
\]

\begin{whiteboard}
\[
(101101)_2
\]

\[
=1(2^5)+0(2^4)+1(2^3)+1(2^2)+0(2)+1
\]

\[
=32+8+4+1
\]

\[
\boxed{(101101)_2=(45)_{10}}
\]
\end{whiteboard}

% ============================================================
\section{Converting Decimal Integers to Another Radix}
% ============================================================

To convert a decimal integer into radix \(r\):

\begin{enumerate}
    \item Divide the number by \(r\).
    \item Record the remainder.
    \item Divide the quotient by \(r\) again.
    \item Continue until the quotient is zero.
    \item Read the remainders from last to first.
\end{enumerate}

Consider the conversion of \(13_{10}\) to binary.

\begin{whiteboard}
\[
\boxed{\text{Decimal integer}\rightarrow\text{base }r}
\]

\[
\text{Repeatedly divide by }r.
\]

\[
\text{Read the remainders upward.}
\]

\[
\begin{array}{rcl}
13/2 &=& 6\quad R1\\
6/2  &=& 3\quad R0\\
3/2  &=& 1\quad R1\\
1/2  &=& 0\quad R1
\end{array}
\]

Reading upward:

\[
\boxed{(13)_{10}=(1101)_2}
\]
\end{whiteboard}

% ============================================================
\section{Converting Decimal Fractions}
% ============================================================

The fractional part of a decimal number is converted differently.

To convert a decimal fraction into radix \(r\):

\begin{enumerate}
    \item Multiply the fraction by \(r\).
    \item Record the integer part of the result.
    \item Multiply the remaining fractional part by \(r\).
    \item Continue until the fraction becomes zero or the desired precision
    has been obtained.
    \item Read the generated integer digits from first to last.
\end{enumerate}

For example, convert

\[
(0.625)_{10}
\]

to binary.

\begin{whiteboard}
\[
\boxed{\text{Decimal fraction}\rightarrow\text{base }r}
\]

\[
\text{Repeatedly multiply by }r.
\]

\[
\text{Read integer parts downward.}
\]

\[
0.625(2)=1.25\Rightarrow1
\]

\[
0.25(2)=0.50\Rightarrow0
\]

\[
0.50(2)=1.00\Rightarrow1
\]

\[
\boxed{(0.625)_{10}=(0.101)_2}
\]
\end{whiteboard}

For a mixed number, the integer and fractional portions are converted
separately and then joined at the radix point.

% ============================================================
\section{Binary, Octal, and Hexadecimal}
% ============================================================

Octal and hexadecimal are especially convenient when working with binary
numbers because their radices are powers of two:

\[
8=2^3,
\qquad
16=2^4.
\]

Therefore,

\[
\boxed{1\text{ octal digit}=3\text{ binary bits}}
\]

and

\[
\boxed{1\text{ hexadecimal digit}=4\text{ binary bits}}.
\]

This makes conversion between these systems a matter of grouping bits rather
than performing arithmetic.

\begin{whiteboard}
\[
8=2^3
\Rightarrow
1\text{ octal digit}=3\text{ bits}
\]

\[
16=2^4
\Rightarrow
1\text{ hex digit}=4\text{ bits}
\]

For hexadecimal:

\[
1010\quad1111\quad0110\quad0011
\]

\[
\phantom{1010}\ A\qquad\ F\qquad\ 6\qquad\ 3
\]

\[
\boxed{(1010111101100011)_2=(AF63)_{16}}
\]

For octal:

\[
\boxed{\text{Group binary bits in groups of 3.}}
\]
\end{whiteboard}

Octal and hexadecimal are convenient ways of displaying the contents of
binary registers because they greatly shorten long strings of \(0\)'s and
\(1\)'s.

% ============================================================
\section{Binary Numbers Versus Binary Codes}
% ============================================================

A crucial distinction must now be made.

A number can be \emph{converted} into binary, or its decimal digits can be
\emph{encoded} with binary patterns. These are not the same operation.

For example,

\[
(99)_{10}=(1100011)_2.
\]

This is the ordinary binary representation of the numerical quantity ninety-nine.

In binary-coded decimal, or BCD, each decimal digit is instead encoded
separately using four bits:

\[
9\rightarrow1001.
\]

Therefore,

\[
99_{\text{BCD}}=1001\ 1001.
\]

\begin{whiteboard}
\[
\boxed{\text{Binary number}\neq\text{binary-coded decimal}}
\]

\[
99_{10}=1100011_2
\]

but

\[
99_{\text{BCD}}=1001\ 1001
\]

\[
\boxed{\text{BCD encodes each decimal digit separately.}}
\]
\end{whiteboard}

The BCD assignments are

\[
\begin{array}{c|c}
\text{Decimal Digit} & \text{BCD} \\ \hline
0&0000\\
1&0001\\
2&0010\\
3&0011\\
4&0100\\
5&0101\\
6&0110\\
7&0111\\
8&1000\\
9&1001
\end{array}
\]

Because four bits permit \(2^4=16\) different patterns but only ten decimal
digits exist, six BCD patterns are unused.

\begin{whiteboard}
\[
0\rightarrow0000,\quad\ldots,\quad9\rightarrow1001
\]

\[
\boxed{1010\text{ through }1111\text{ are unused in BCD.}}
\]
\end{whiteboard}

% ============================================================
\section{Alphanumeric Representation and ASCII}
% ============================================================

Computers must also represent letters, digits, punctuation marks, and special
characters.

An \emph{alphanumeric code} assigns a binary pattern to each character in a
character set.

The standard code discussed by Mano is ASCII, the American Standard Code
for Information Interchange.

ASCII uses seven bits, permitting

\[
2^7=128
\]

different character codes.

Examples include

\[
A=1000001
\]

and

\[
0=0110000.
\]

\begin{whiteboard}
\[
\boxed{\text{ASCII}}
\]

\[
7\text{ bits}\Rightarrow2^7=128\text{ possible codes}
\]

\[
A=1000001
\]

\[
0=0110000,\qquad9=0111001
\]
\end{whiteboard}

An interesting relationship exists between the ASCII codes for decimal digits
and BCD. The ASCII digits begin with the three high-order bits \(011\).

Thus,

\[
011\,0000
\]

is ASCII \(0\), while

\[
011\,1001
\]

is ASCII \(9\).

Removing the high-order \(011\) leaves the corresponding four-bit BCD digit.

\begin{whiteboard}
\[
\text{ASCII decimal digit: }011\,xxxx
\]

\[
\text{Remove }011
\]

\[
\Downarrow
\]

\[
\boxed{\text{BCD digit }xxxx}
\]
\end{whiteboard}

% ============================================================
\section{Complements}
% ============================================================

Complements are useful because they allow subtraction to be performed through
addition.

For a base-\(r\), \(n\)-digit number \(N\), Mano defines two complements.

The \((r-1)\)'s complement is

\[
(r^n-1)-N.
\]

The \(r\)'s complement is

\[
r^n-N
\]

for nonzero \(N\).

\begin{whiteboard}
\[
\boxed{(r-1)\text{'s complement}=(r^n-1)-N}
\]

\[
\boxed{r\text{'s complement}=r^n-N}
\]

\[
\begin{array}{c|c|c}
 & r-1 & r \\ \hline
\text{decimal} & 9\text{'s} & 10\text{'s}\\
\text{binary} & 1\text{'s} & 2\text{'s}
\end{array}
\]
\end{whiteboard}

% ------------------------------------------------------------
\subsection{1's Complement}
% ------------------------------------------------------------

The 1's complement of a binary number is obtained by changing every

\[
0\rightarrow1
\]

and every

\[
1\rightarrow0.
\]

For example,

\[
1011001
\]

becomes

\[
0100110.
\]

\begin{whiteboard}
\[
\boxed{\text{1's complement: flip every bit}}
\]

\[
1011001
\]

\[
\Downarrow
\]

\[
0100110
\]
\end{whiteboard}

% ------------------------------------------------------------
\subsection{2's Complement}
% ------------------------------------------------------------

The 2's complement is obtained by first taking the 1's complement and then
adding one.

For example,

\[
101100
\]

has 1's complement

\[
010011.
\]

Adding one gives

\[
010011+1=010100.
\]

\begin{whiteboard}
\[
\boxed{\text{2's complement: flip every bit, then add 1}}
\]

\[
101100
\]

\[
\Downarrow
\]

\[
010011
\]

\[
+1
\]

\[
\boxed{010100}
\]
\end{whiteboard}

A useful shortcut is to begin at the right side of the number.

Leave all trailing \(0\)'s and the first \(1\) unchanged. Complement every
remaining bit to the left.

\begin{whiteboard}
\[
\boxed{
\begin{array}{c}
\text{2's-complement shortcut:}\\
\text{from the right, keep trailing zeros and the first 1;}\\
\text{flip everything to the left.}
\end{array}}
\]
\end{whiteboard}

% ============================================================
\section{Subtraction Using Complements}
% ============================================================

Suppose two \(n\)-bit unsigned numbers \(M\) and \(N\) are to be subtracted.

Instead of directly computing

\[
M-N,
\]

we may add the 2's complement of \(N\):

\[
M+(2^n-N).
\]

Algebraically,

\[
M+(2^n-N)=M-N+2^n.
\]

If an end carry occurs, the \(2^n\) term appears as that carry and is discarded.

\begin{whiteboard}
\[
\boxed{M-N=M+(2^n-N)}
\]

\[
M+(2^n-N)=M-N+2^n
\]

\[
\boxed{\text{Subtraction can be implemented with addition hardware.}}
\]
\end{whiteboard}

Consider \(7-5\) using four bits.

\[
7=0111,
\qquad
5=0101.
\]

The 2's complement of \(0101\) is

\[
1011.
\]

Therefore,

\[
0111+1011=1\,0010.
\]

The end carry is discarded, leaving

\[
0010=2.
\]

\begin{whiteboard}
\[
7-5
\]

\[
7=0111,\qquad5=0101
\]

\[
\text{2's comp. of }0101=1011
\]

\[
\begin{array}{r}
0111\\
+1011\\ \hline
1\,0010
\end{array}
\]

\[
\text{discard carry}
\]

\[
\boxed{0010=2}
\]
\end{whiteboard}

% ============================================================
\section{Fixed-Point Representation}
% ============================================================

A binary number may represent an integer, a fraction, or a mixed quantity.

The binary point itself does not ordinarily need to occupy a physical storage
position. Instead, its location can be assumed.

In \emph{fixed-point representation}, the binary point always occupies an
assumed fixed position.

Two common arrangements are:

\begin{itemize}
    \item binary point assumed at the extreme right, producing integers;
    \item binary point assumed at the extreme left, producing fractions.
\end{itemize}

\begin{whiteboard}
\[
\boxed{\text{Fixed point: the radix-point position is assumed.}}
\]

\[
\text{point at right}\Rightarrow\text{integer}
\]

\[
\text{point at left}\Rightarrow\text{fraction}
\]
\end{whiteboard}

% ============================================================
\section{Signed Integer Representation}
% ============================================================

A sign must also be represented with a bit.

The convention used is

\[
0=\text{positive},
\qquad
1=\text{negative}.
\]

For negative integers, Mano discusses three representations:

\begin{enumerate}
    \item signed magnitude,
    \item signed 1's complement,
    \item signed 2's complement.
\end{enumerate}

Consider \(14\) in an 8-bit register.

\[
+14=00001110.
\]

The three representations of \(-14\) are:

\[
\begin{array}{c|c}
\text{Representation} & -14 \\ \hline
\text{Signed magnitude} & 10001110\\
\text{Signed 1's complement} & 11110001\\
\text{Signed 2's complement} & 11110010
\end{array}
\]

\begin{whiteboard}
\[
0=\text{positive sign},
\qquad
1=\text{negative sign}
\]

\[
+14=00001110
\]

\[
\begin{array}{c|c}
\text{Representation of }-14 & \text{Bits}\\ \hline
\text{signed magnitude} & 10001110\\
\text{1's complement} & 11110001\\
\text{2's complement} & 11110010
\end{array}
\]
\end{whiteboard}

Signed magnitude resembles ordinary written arithmetic but makes computer
arithmetic comparatively awkward.

The 1's-complement representation also has the inconvenience of possessing
two representations of zero:

\[
+0
\qquad\text{and}\qquad
-0.
\]

For this reason, Mano's discussion of signed binary arithmetic concentrates
on the 2's-complement system.

\begin{whiteboard}
\[
\text{signed magnitude: awkward arithmetic}
\]

\[
\text{1's complement: }+0\text{ and }-0
\]

\[
\boxed{\text{2's complement is preferred for signed arithmetic.}}
\]
\end{whiteboard}

% ============================================================
\section{Signed 2's-Complement Addition}
% ============================================================

The great advantage of 2's complement is that the arithmetic rule is extremely
simple:

\begin{enumerate}
    \item Add both numbers, including their sign bits.
    \item Discard any carry out of the sign-bit position.
\end{enumerate}

Negative results remain naturally expressed in 2's-complement form.

\begin{whiteboard}
\[
\boxed{\text{2's-complement addition}}
\]

\[
\text{Add all bits, including sign bits.}
\]

\[
\text{Discard carry out of the sign bit.}
\]
\end{whiteboard}

Consider

\[
(+6)+(-13).
\]

In eight bits,

\[
+6=00000110.
\]

To represent \(-13\),

\[
13=00001101,
\]

so

\[
-13=11110011.
\]

Therefore,

\[
\begin{array}{r}
00000110\\
+11110011\\ \hline
11111001
\end{array}
\]

The result has sign bit \(1\), so it is negative.

Taking its 2's complement,

\[
11111001
\rightarrow
00000111.
\]

Thus,

\[
11111001=-7.
\]

\begin{whiteboard}
\[
+6=00000110
\]

\[
-13=11110011
\]

\[
\begin{array}{r}
00000110\\
+11110011\\ \hline
11111001
\end{array}
\]

\[
11111001\xrightarrow{\text{2's comp.}}00000111
\]

\[
\boxed{11111001=-7}
\]
\end{whiteboard}

% ============================================================
\section{Signed 2's-Complement Subtraction}
% ============================================================

Subtraction follows directly from

\[
A-B=A+(-B).
\]

Therefore, to subtract \(B\) from \(A\):

\begin{enumerate}
    \item Take the 2's complement of \(B\).
    \item Add it to \(A\).
    \item Discard any carry out of the sign position.
\end{enumerate}

\begin{whiteboard}
\[
\boxed{A-B=A+(-B)}
\]

\[
\boxed{\text{2's-complement the subtrahend, then add.}}
\]
\end{whiteboard}

This means that the same basic binary addition circuitry can be used for both
addition and subtraction.

% ============================================================
\section{Overflow}
% ============================================================

A computer register has finite width. If the true arithmetic result cannot fit
within that width, an \emph{overflow} occurs.

The interpretation of overflow depends upon whether the operands are
unsigned or signed.

% ------------------------------------------------------------
\subsection{Unsigned Overflow}
% ------------------------------------------------------------

For unsigned numbers, overflow is indicated by a carry out of the most
significant bit.

\begin{whiteboard}
\[
\boxed{\text{Unsigned overflow: carry out of the MSB}}
\]
\end{whiteboard}

% ------------------------------------------------------------
\subsection{Signed 2's-Complement Overflow}
% ------------------------------------------------------------

For signed 2's-complement numbers, the final carry alone does not determine
overflow.

Overflow can occur when two operands have the same sign but the stored result
has the opposite sign.

Thus:

\[
(+)+(+)\rightarrow(-)
\]

indicates overflow, and

\[
(-)+(-)\rightarrow(+)
\]

also indicates overflow.

\begin{whiteboard}
\[
\boxed{
\text{same-sign operands + opposite-sign result}
\Rightarrow
\text{overflow}}
\]
\end{whiteboard}

Mano gives a hardware test based upon the carry into and carry out of the sign
bit.

Let

\[
C_{\text{in}}
\]

be the carry into the sign position and

\[
C_{\text{out}}
\]

be the carry out of it.

Then

\[
V=C_{\text{in}}\oplus C_{\text{out}}.
\]

Therefore,

\[
C_{\text{in}}\neq C_{\text{out}}
\]

indicates overflow.

\begin{whiteboard}
\[
\boxed{V=C_{\text{in sign}}\oplus C_{\text{out sign}}}
\]

\[
C_{\text{in}}\neq C_{\text{out}}
\Rightarrow
V=1
\]

\[
\boxed{\text{Overflow}}
\]
\end{whiteboard}

An \(n\)-bit 2's-complement register represents values in the range

\[
-2^{n-1}
\leq x \leq
2^{n-1}-1.
\]

For eight bits,

\[
-128\leq x\leq127.
\]

Thus, although

\[
70+80=150,
\]

the value \(150\) cannot be represented as an 8-bit signed 2's-complement
integer.

\begin{whiteboard}
\[
\boxed{-2^{n-1}\leq x\leq2^{n-1}-1}
\]

For \(n=8\):

\[
-128\leq x\leq127
\]

\[
70+80=150>127
\]

\[
\boxed{\text{overflow}}
\]
\end{whiteboard}

% ============================================================
\section{Decimal Fixed-Point Representation}
% ============================================================

Decimal values can also be stored and manipulated directly in a binary-coded
form.

Using BCD requires four bits for every decimal digit.

For example,

\[
4385
\]

is represented as

\[
0100\ 0011\ 1000\ 0101.
\]

\begin{whiteboard}
\[
\boxed{\text{Decimal fixed point uses coded decimal digits.}}
\]

\[
4385_{\text{BCD}}
=
0100\ 0011\ 1000\ 0101
\]

\[
\boxed{4\text{ bits per decimal digit}}
\]
\end{whiteboard}

Decimal representation requires more storage than the equivalent pure binary
number, and decimal arithmetic circuitry is more complicated.

Its advantage is that decimal-oriented input, output, and processing can be
performed without repeatedly converting between decimal and binary forms.

\begin{whiteboard}
\[
\text{BCD: more storage}
\]

\[
\text{BCD arithmetic: more complex hardware}
\]

\[
\text{advantage: convenient for decimal I/O and processing}
\]
\end{whiteboard}

The decimal counterpart of binary 2's-complement arithmetic is the
10's-complement system.

\begin{whiteboard}
\[
\boxed{\text{Decimal analogue of 2's complement: 10's complement}}
\]
\end{whiteboard}

% ============================================================
\section{Floating-Point Representation}
% ============================================================

Fixed-point representation assumes that the radix point remains in one fixed
position.

Floating-point representation instead stores information that allows the
effective radix-point position to vary.

Mano expresses a floating-point quantity as

\[
m\times r^e,
\]

where

\[
m=\text{mantissa}
\]

and

\[
e=\text{exponent}.
\]

The radix \(r\) and the location of the radix point within the mantissa are
understood rather than explicitly stored.

\begin{whiteboard}
\[
\boxed{\text{Floating point}=m\times r^e}
\]

\[
m=\text{mantissa}
\]

\[
e=\text{exponent}
\]
\end{whiteboard}

For example,

\[
6132.789
=
0.6132789\times10^4.
\]

\begin{whiteboard}
\[
6132.789
=
0.6132789\times10^4
\]
\end{whiteboard}

The same principle applies in binary.

For example,

\[
1001.11_2
=
0.100111_2\times2^4.
\]

\begin{whiteboard}
\[
1001.11_2
=
0.100111_2\times2^4
\]
\end{whiteboard}

% ------------------------------------------------------------
\subsection{Normalization}
% ------------------------------------------------------------

A floating-point number is \emph{normalized} when the most significant digit
of the mantissa is nonzero.

Normalization moves the radix point while changing the exponent by a
corresponding amount so that the numerical value remains unchanged.

For example,

\[
0.001101\times2^5
\]

may be normalized by shifting the mantissa two positions left:

\[
0.1101\times2^3.
\]

The value has not changed.

\begin{whiteboard}
\[
\boxed{\text{normalized: most significant mantissa digit is nonzero}}
\]

\[
0.001101\times2^5
\]

shift left two places:

\[
0.1101\times2^3
\]

\[
\boxed{\text{same value}}
\]

\[
\boxed{\text{Normalization maximizes available precision.}}
\]
\end{whiteboard}

Floating-point arithmetic requires more complicated hardware than fixed-point
arithmetic, but it permits a much wider range of numerical magnitudes.

% ============================================================
\section{Gray Code}
% ============================================================

The reflected binary code, usually called \emph{Gray code}, has the important
property that adjacent values differ in only one bit.

A portion of the Gray-code sequence is

\[
000,\ 001,\ 011,\ 010,\ 110,\ 111,\ldots
\]

\begin{whiteboard}
\[
\boxed{\text{Gray code: adjacent values differ by one bit.}}
\]

\[
000\rightarrow001\rightarrow011\rightarrow010
\]

\[
\rightarrow110\rightarrow111\rightarrow\cdots
\]

\[
\boxed{\text{Only one bit changes at each step.}}
\]
\end{whiteboard}

This property is useful when a physical or digital state passes from one value
to the next because it reduces ambiguity during the transition.

\begin{whiteboard}
\[
\boxed{\text{Advantage: reduced ambiguity during transitions}}
\]
\end{whiteboard}

% ============================================================
\section{Other Decimal Codes}
% ============================================================

BCD is not the only possible binary code for decimal digits.

Mano also discusses:

\begin{itemize}
    \item 8421 BCD,
    \item 2421 code,
    \item Excess-3 code,
    \item Excess-3 Gray code.
\end{itemize}

BCD may be viewed as an \(8421\) weighted code because the bit positions have
weights

\[
8,\ 4,\ 2,\ 1.
\]

The 2421 code is another weighted code.

The Excess-3 code is obtained by adding binary \(3\), or \(0011\), to the BCD
representation of each decimal digit.

The 2421 and Excess-3 codes have a self-complementing property: complementing
all bits produces the code representation of the 9's complement of the
decimal digit.

\begin{whiteboard}
\[
\boxed{\text{Other decimal codes}}
\]

\[
\text{BCD}=8421\text{ weighted code}
\]

\[
\text{2421: weighted and self-complementing}
\]

\[
\text{Excess-3}=\text{BCD}+0011
\]

\[
\boxed{
\text{2421 and Excess-3: bit complement}
\Rightarrow
\text{9's complement}}
\]
\end{whiteboard}

% ============================================================
\section{Other Alphanumeric Codes}
% ============================================================

Mano also identifies other character-code arrangements.

ASCII uses seven bits for each character and commonly receives an additional
parity bit during transmission.

EBCDIC, associated with IBM equipment, uses eight bits per character.

Some systems may also use shorter internal character codes when a smaller
character set is sufficient.

\begin{whiteboard}
\[
\boxed{\text{Alphanumeric codes}}
\]

\[
\text{ASCII: 7-bit character code}
\]

\[
\text{EBCDIC: 8-bit character code}
\]

\[
\text{smaller internal codes may also be used}
\]
\end{whiteboard}

% ============================================================
\section{Error Detection and Parity}
% ============================================================

Binary information sent through a communication medium may be affected by
noise.

A transmitted

\[
1
\]

may become

\[
0,
\]

or a transmitted

\[
0
\]

may become

\[
1.
\]

An error-detection code adds information that allows the receiver to recognize
that some transmission errors have occurred.

The most common method discussed by Mano is the \emph{parity bit}.

\begin{whiteboard}
\[
\boxed{\text{Transmission error: }1\leftrightarrow0}
\]

\[
\boxed{\text{Parity bit: extra bit used for error detection}}
\]
\end{whiteboard}

% ------------------------------------------------------------
\subsection{Even and Odd Parity}
% ------------------------------------------------------------

With \emph{even parity}, the parity bit is chosen so that the total number of
\(1\)'s in the transmitted word is even.

With \emph{odd parity}, the parity bit is chosen so that the total number of
\(1\)'s is odd.

\begin{whiteboard}
\[
\boxed{\text{Even parity: total number of 1s is even}}
\]

\[
\boxed{\text{Odd parity: total number of 1s is odd}}
\]
\end{whiteboard}

Suppose the message is

\[
1011.
\]

It contains three \(1\)'s.

To obtain even parity, the parity bit must therefore be

\[
P=1,
\]

giving four \(1\)'s in the transmitted word.

\begin{whiteboard}
Message:

\[
1011
\]

Number of \(1\)'s:

\[
3
\]

For even parity:

\[
P=1
\]

Total number of \(1\)'s:

\[
4
\]

\[
\boxed{\text{even parity satisfied}}
\]
\end{whiteboard}

% ------------------------------------------------------------
\subsection{Parity and Exclusive-OR}
% ------------------------------------------------------------

Parity generation connects directly to the exclusive-OR operation studied in
digital logic.

For a three-bit message \(x,y,z\),

\[
P_{\text{even}}
=
x\oplus y\oplus z.
\]

The exclusive-OR function is \(1\) whenever an odd number of its inputs are
\(1\).

Odd parity therefore requires the complement:

\[
P_{\text{odd}}
=
\overline{x\oplus y\oplus z}.
\]

\begin{whiteboard}
\[
\boxed{P_{\text{even}}=x\oplus y\oplus z}
\]

\[
\boxed{
P_{\text{odd}}
=
\overline{x\oplus y\oplus z}}
\]

\[
\text{XOR is an odd function.}
\]
\end{whiteboard}

% ------------------------------------------------------------
\subsection{Parity Generator and Checker}
% ------------------------------------------------------------

At the sending end, a parity generator calculates the required parity bit.

The message and parity bit are then transmitted.

At the receiving end, a parity checker determines whether the received bit
pattern still has the required parity.

\begin{whiteboard}
\[
\text{message}
\rightarrow
\boxed{\text{parity generator}}
\rightarrow
\text{channel}
\rightarrow
\boxed{\text{parity checker}}
\rightarrow
\text{error indication}
\]
\end{whiteboard}

% ------------------------------------------------------------
\subsection{What Parity Can and Cannot Detect}
% ------------------------------------------------------------

A single parity bit detects an odd number of changed bits.

Thus it detects:

\[
1,\ 3,\ 5,\ldots
\]

bit errors.

However, if an even number of bits change, the parity may remain correct and
the error can go undetected.

\begin{whiteboard}
\[
\boxed{\text{Parity detects an odd number of bit errors.}}
\]

\[
1,\ 3,\ 5,\ldots
\]

\[
\boxed{\text{An even number of bit errors may go undetected.}}
\]
\end{whiteboard}

Parity therefore provides \emph{error detection}, not complete error protection
and not automatic error correction.

% ============================================================
\section{Chapter 3 Summary}
% ============================================================

The fundamental theme of this chapter is that computers physically store
binary information, but those bits acquire meaning only through an agreed
representation.

We have seen that bits can represent:

\begin{itemize}
    \item binary, octal, decimal, and hexadecimal quantities,
    \item decimal digits through BCD,
    \item characters through codes such as ASCII,
    \item positive and negative integers,
    \item fixed-point and floating-point quantities,
    \item alternate codes such as Gray code,
    \item and redundant information used for error detection.
\end{itemize}

The particularly important architectural consequence is that a representation
is chosen not merely for human readability but also because it can simplify
the hardware required to manipulate the information.

Two's-complement representation is a prime example: subtraction can be
reduced to addition, allowing common arithmetic hardware to serve both
operations.

\begin{whiteboard}
\[
\boxed{
\text{bits}
\rightarrow
\text{number systems}
\rightarrow
\text{codes}
\rightarrow
\text{signed arithmetic}
}
\]

\[
\boxed{
\text{fixed/floating point}
\rightarrow
\text{error detection}
}
\]

\[
\boxed{\text{Representation affects both meaning and hardware.}}
\]
\end{whiteboard}

\end{document}